% fig_7_39.tex — 7-54: 竖直圆环内放直角三角板（斜边2R，短边R）
\begin{tikzpicture}[font=\small,>=stealth]
  % 圆环
  \draw[thick] (0,0) circle (2);
  % 圆心 O
  \fill (0,0) circle (1.5pt);
  \node[below right] at (0,0) {$O$};
  % 铅垂线
  \draw[dashed,gray] (0,2.3) -- (0,-2.3);
  % 直角三角板（斜边=2R，短边=R，放在环内底部平衡）
  % 斜边 = 2R = 环的直径的一端到另一端？不，斜边=2R=环直径
  % 短边=R，另一条直角边=√(4R²-R²)=R√3
  % 平衡位置：斜边在底部
  % 三角板的三个顶点在环上
  % 设斜边水平放在底部，两端在环上
  % 短边=R 从一端向上
  % 顶点1: (-R, -√(R²-R²)) = (-R, 0) → 不对
  % 用参数：斜边长2R，放在圆环内
  % 圆环半径=R，三角板斜边=2R=直径
  % 所以斜边两端在圆环直径两端
  % 斜边水平时：(-R,0) 和 (R,0)
  % 短边R从(-R,0)出发，垂直方向？
  % 直角在短边和长直角边的交点
  % 短边=R，从(-R,0)向上R → (-R,R)
  % 但(-R,R)到(R,0)的距离 = √(4R²+R²) = R√5 ≠ R√3
  % 重新考虑：直角三角形斜边2R，一条直角边R
  % 直角顶点C，两条直角边CA=R, CB=R√3, 斜边AB=2R
  % 放在环内：A和B在环上，C在环内
  % 平衡：重心在最低点
  % 三角板顶点坐标（斜边AB水平在底部）
  % A=(-R,0), B=(R,0), C在AB上方
  % C到AB的距离h = R*R√3/(2R) = R√3/2
  % 但C的位置需要满足CA=R, CB=R√3
  % 从A=(-R,0)出发，CA=R：C在以A为圆心R为半径的圆上
  % 从B=(R,0)出发，CB=R√3：C在以B为圆心R√3为半径的圆上
  % C = (-R + R*cosα, R*sinα) 其中α满足
  % (-R+R*cosα-R)² + (R*sinα)² = 3R²
  % (R*cosα-2R)² + R²sin²α = 3R²
  % R²cos²α - 4R²cosα + 4R² + R²sin²α = 3R²
  % R² - 4R²cosα + 4R² = 3R²
  % 5R² - 4R²cosα = 3R²
  % cosα = 1/2 → α = 60°
  % C = (-R + R/2, R√3/2) = (-R/2, R√3/2)
  \draw[thick,fill=orange!15] (-2,0) -- (2,0) -- ({-2*0.5},{2*0.866}) -- cycle;
  \node[font=\footnotesize] at (-0.7,0.5) {三角板};
  % 标注顶点
  \fill (-2,0) circle (1.5pt) node[left] {$A$};
  \fill (2,0) circle (1.5pt) node[right] {$B$};
  \fill ({-2*0.5},{2*0.866}) circle (1.5pt) node[above] {$C$};
  % 标注边长
  \node[below,font=\footnotesize] at (0,-0.2) {斜边 $2R$};
  \node[left,font=\footnotesize] at (-1.6,0.5) {$R$};
  \node[right,font=\footnotesize] at (0.6,0.5) {$R\sqrt{3}$};
\end{tikzpicture}
